LLM benchmarks are the primary tool for comparing language models, guiding deployment, funding, and ranking decisions. Every evaluation involves design choices (prompt template, random seed, numerical precision, answer ordering) typically fixed to single values and assumed inconsequential. Growing evidence challenges this assumption: prompt formatting can shift accuracy by 76 percentage points, and answer ordering can flip model rankings. Yet each study examines one factor in isolation, leaving unanswered which sources dominate when all vary simultaneously and how many items suffice. We address this with Generalizability theory, a psychometric framework for multi-source variance decomposition. In a factorial experiment crossing seven facets across eight models and five benchmarks (2.3M records), we find that item difficulty is the largest variance component (34--38\% in multiple-choice), and the model$\times$item interaction is the second largest (30--36\% in MC, 15--17\% in free-form), exceeding every manipulated design factor by an order of magnitude---models disagree on which items are difficult, not in overall ability. A 3-prompt $\times$ 3-seed protocol absorbs this interaction, reducing items from 552 to 178 ($G \geq 0.80$). We observe a format-associated pattern ($n = 5$, exploratory): item-related variance concentrates in multiple-choice (68--72\%) but disperses in free-form tasks (31--41\%). We release an open-source G-theory toolkit\footnote{Code and data are available at \url{https://anonymous.4open.science/r/nlp_variance_gtheory-11E1/}.}.
